\chapter*{\chapterstyle{II --- Espaces Topologiques}}
\addcontentsline{toc}{section}{Espaces Topologiques}

Soit \(E\) un ensemble \textbf{non-vide} et \(\mathcal{T}\) une ensemble de parties de \(E\) (qu'on appelera \textbf{ouverts}) qui vérifient les propriétés suivantes:
\begin{align*}
   &\bullet \;\;\; \text{\(E\) et \(\emptyset\) sont des ouverts} \\
   &\bullet \;\; \text{Toute union \textbf{quelconque} d'ouverts est un ouvert} \\
   &\bullet \;\; \text{Toute intersection \textbf{finie} d'ouverts est un ouvert}
\end{align*}
Alors le couple \((E, \mathcal{T})\) est appellé \textbf{espace topologique} et \(\mathcal{T}\) est appellée \textbf{topologie} sur \(E\).\<

On dit alors qu'une partie de \(E\) est \textbf{fermée} si son complémentaire est \textbf{ouvert}.
\subsection*{\subsecstyle{Intérieur {:}}}
Soit \(A \subseteq E\), alors on dit que \(x \in A\) est un point \textbf{intérieur} à \(A\) si et seulement si:
\customBox{width=10cm}{
   Il existe un ouvert \(\mathcal{O}\) qui contient \(x\) tel que \(\mathcal{O} \subseteq A\).
}

On peut alors définir \textbf{l'intérieur} d'une partie de \(E\) comme étant l'ensemble de ses points intérieurs. C'est une application de l'ensemble des parties de \(E\) dans lui-meme qui est \textbf{idempotente, croissante et anti-extensive}\footnote[2]{On appelle anti-extensive une application telle que \(\phi(S) \subseteq S\) (non-usuel)}.\<

On peut alors caractériser l'intérieur de \(A\) comme étant \textbf{le plus grand ouvert contenu dans \(A\)}. En particulier, c'est l'union de tout les ouverts contenus dans \(A\) et on a:
\[
    \text{int}(A) := \bigcup_{\mathcal{O} \subseteq{A}} \mathcal{O}
\]

Enfin, on peut caractériser les ouverts par leur intérieur, en effet \(\mathcal{O}\) est un ouvert si et seulement si:
\customBox{width=3cm}{
   \(\mathcal{O} = \text{int}(\mathcal{O})\)
}
\subsection*{\subsecstyle{Adhérence {:}}}
Soit \(A \subseteq E\), alors on dit que \(x \in A\) est un point \textbf{adhérent} à \(A\) si et seulement si:
\customBox{width=10cm}{
   Pour tout ouvert \(\mathcal{O}\) qui contient \(x\), alors \(\mathcal{O} \cap A \neq \emptyset\).
}
On peut alors définir \textbf{l'adhérence} d'une partie de \(E\) comme étant l'ensemble de ses points adhérents. C'est une application de l'ensemble des parties de \(E\) dans lui-meme qui est \textbf{idempotente, croissante, et extensive}\footnote[3]{On appelle extensive une application telle que \( S \subseteq \phi(S) \)}.\<

On peut alors caractériser l'adhérence de \(A\) comme étant \textbf{le plus petit fermé qui contient \(A\)}. En particulier, c'est l'intersection de tout les fermés qui contiennent \(A\) et on a:
\[
    \text{adh}(A) := \bigcap_{\mathcal{F} \supseteq {A}} \mathcal{F}
\]

Enfin, on peut caractériser les fermés par leur adhérence, en effet \(\mathcal{F}\) est un fermé si et seulement si:
\customBox{width=3cm}{
   \(\mathcal{F} = \text{adh}(\mathcal{F})\)
}
\subsection*{\subsecstyle{Frontière {:}}}
On appelle \textbf{frontière} d'une partie \(A\) de \(E\) et on note \(\partial A\) la partie qui vérifie:
\customBox{width=5cm}{
   \(\partial A = \text{adh}{A} \backslash \text{int}{A}\)
}
\subsection*{\subsecstyle{Propriétés {:}}}
Soit \(A \subseteq E\), on peut alors montrer plusieurs propriétés de l'intérieur et de l'adhérence:
\customBox{width=8cm}{
   \( \text{int}(A)^c = \text{adh}(A^c)
   \quad\quad\quad\quad
   \text{adh}(A)^c = \text{int}(A^c)
   \)
}  
On peut aussi montrer les différentes relations suivantes\footnote[1]{Pour comprendre les inclusions strictes, considérer \(\Q\) et \(\R \backslash \Q\) (en tant que parties de \(\R\)).} vis à vis des opérations ensemblistes:
\customBox{width=12cm}{
   \begin{align*}
      \text{int}(A \cap B) = \text{int}(A) \cap \text{int}(B) \quad\quad\quad\quad \text{adh}(A \cup B) = \text{adh}(A) \cup \text{adh}(B) \\
      \text{int}(A \cup B) \supseteq \text{int}(A) \cup \text{int}(B) \quad\quad\quad\quad \text{adh}(A \cap B) \subseteq \text{adh}(A) \cap \text{adh}(B)
   \end{align*}
}  
\subsection*{\subsecstyle{Topologies usuelles {:}}}
On peut distinguer plusieurs topologies usuelles:\<

La \textbf{topologie grossière}, alors les ouverts sont seulements \(E\) et \(\emptyset\).\+
La \textbf{topologie discrète}, alors les ouverts sont toutes les parties de \(E\).\+
La \textbf{topologie de l'ordre} sur un ensemble ordonné, alors les ouverts sont les intervalles pour la relation\footnote[2]{C'est la topologie usuelle de \(\R\).}. 
\subsection*{\subsecstyle{Topologie induite {:}}}
Soit \(A \subseteq E\), alors \(E\) induit une topologie naturelle sur \(A\) appellée \textbf{topologie induite} définie par:
\[
   \mathcal{T}_A := \Bigl\{ A \cap \mathcal{O} \; ; \; \mathcal{O} \subseteq \mathcal{T}_E \Bigl\}  
\]
\begin{center}
   \textit{Un ouvert de \(A\) pour la topologie induite est l'intersection d'un ouvert de l'espace total avec \(A\)}
\end{center}
En particulier on peut facilement représenter cette topologie dans l'espace métrique \(\R^2\):
\begin{center}
   \begin{tikzpicture}[
      >=stealth,scale=1,line cap=round,
      bullet/.style={circle,inner sep=0.5pt,fill}
   ]
      \filldraw[color=BrightBlue1!60, fill=BrightBlue1!5, dashed] (4.25,2.75) circle (2) node {$\mathcal{O} \in \mathcal{T}_E$};
      \draw[->] (-1,0) -- (9,0);
      \draw[->] (0,-1) -- (0, 5.5);
      \draw node at (8.25,5) {$E$};
      \draw[color=DarkBlue1, domain=0:8, smooth]   plot (\x,{sin(\x r) + 1.5})   node[right] {$A$}; 
      \draw[color=BrightRed1, domain=2.35:3.92, smooth, line width=0.5mm]   plot (\x,{sin(\x r) + 1.5})   node[below left] {$\mathcal{O} \in \mathcal{T}_A$}; 
   \end{tikzpicture}   
\end{center}

\chapter*{\chapterstyle{II --- Espaces Métriques}}
\addcontentsline{toc}{section}{Espaces Métriques}
Soit \(E\) un ensemble, le cas le plus courant d'espace topologique est le cas des espaces munis d'une application apellée \textbf{distance} ou aussi métrique, c'est une application de \(E \times E\) dans \(\R^+\) qui vérifie:
\begin{align*}
   &\bullet \;\; \text{Elle est \textbf{symétrique}.} \\
   &\bullet \;\; \text{Elle vérifie \textbf{l'inégalité triangulaire}.} \\
   &\bullet \;\; \text{Deux points sont égaux si et seulement si la distance entre ces points est nulle.}
\end{align*}
Alors, l'espace \(E\) muni de cette distance notée \(d\) est alors appelé \textbf{espace métrique}, et on le note \((E, d)\).\<

Si on considére le cas particulier des \textbf{espaces vectoriels normés} (et donc a fortiori des espaces euclidiens), alors la norme induit une distance sur \(E\) et en fait un espace métrique.

\subsection*{\subsecstyle{Topologie {:}}}
On appelle \textbf{boule ouverte} de centre \(x\) et de rayon \(r > 0\), la partie:
\customBox{width=6cm}{
   \(\mathcal{B}(x, r) := \Bigl\{ y \in E \; ; \; d(x, y) < r\Bigl\}\)
}
On peut par ailleurs définir \textbf{la boule fermée} de centre \(x\) et de rayon \(r > 0\) comme étant la partie:
\customBox{width=6cm}{
   \(\mathcal{B}[x, r] := \Bigl\{ y \in E \; ; \; d(x, y) \leq r\Bigl\}\)
}
La distance sur \(E\) induit alors naturellement une topologie sur \(E\) dont les ouverts sont \textbf{les boules ouvertes} (et les unions quelconques de boules ouvertes), la notion d'intérieur, d'adhérence, et de frontière s'étends alors aisément à ce cas particulier d'ouvert.


\subsection*{\subsecstyle{Propriétés {:}}}
On peut alors montrer que dans un espace métrique on a:
\customBox{width=9cm}{
   \( \text{int}(\mathcal{B}[x, r]) \supseteq \mathcal{B}(x, r)
   \quad\quad\quad\quad
   \text{adh}(\mathcal{B}(x, r)) \subseteq \mathcal{B}[x, r]
   \)
}
En particulier, dans le cas d'un \textbf{espace vectoriel normé}, l'homogénéité de la norme nous permet d'avoir l'égalité plutot que l'inclusion pour ces deux propriétés.

\subsection*{\subsecstyle{Exemples de boules en dimension finie {:}}}
On considère ici l'espace vectoriel usuel \(\R^2\), on peut alors définir les \textbf{normes de Hölder} définies comme suit:
\customBox{width=6cm}{
   \({||u||}_p = (|u_1|^p + |u_2|^p)^{\frac{1}{p}}\)
}
En particulier pour \(p=2\), on reconnait \textbf{la norme euclidienne standard}. On peut aussi définir \textbf{la norme infini} par:
\customBox{width=6cm}{
   \({||u||}_\infty = \sup\{|u_1|, |u_2|\}\)
}
\pagebreak

On peut représenter les boules pour quelques valeurs de \(p\):
\begin{center}
   \begin{tikzpicture}[
      >=stealth,scale=1,line cap=round,
      bullet/.style={circle,inner sep=0.5pt,fill}
   ]
      \draw[->] (-2,0) -- (2,0);
      \draw[->] (0,-2) -- (0, 2);

      \draw[DarkBlue1] (-1.75,0)--(-1.75,1.75)--(0,1.75)--(1.75,1.75)--(1.75,0)--(1.75,-1.75)--(0,-1.75)--(-1.75,-1.75)--cycle;
      \draw[BrightBlue1] (-1.75,0)--(0,1.75)--(1.75,0)--(0,-1.75)--cycle;
      \draw[BrightRed1](0,0) circle (1.75);

      \draw[DarkBlue1] node at (2.5, 1.60) {$p = \infty$};
      \draw[BrightRed1] node at (2.4, 1.20) {$p = 2$};
      \draw[BrightBlue1] node at (2.4, 0.75) {$p = 1$};
   \end{tikzpicture}   
\end{center}

\subsection*{\subsecstyle{Exemples de boules en dimension infinie {:}}}
On considère ici l'espace vectoriel des \textbf{fonction continues} sur l'intervalle \(\icc{a}{b}\), on peut alors étendre les \textbf{normes de Hölder} définies comme suit:
\customBox{width=6cm}{
   \({||f||}_p = (\int_{a}^{b} |f(t)|^p \d t)^{\frac{1}{p}}\)
}
En particulier pour \(p=2\), c'est \textbf{la norme euclidienne standard}. On peut aussi définir \textbf{la norme infini} (aussi appelée \textbf{norme de la convergence uniforme}) par:
\customBox{width=6cm}{
   \({||f||}_\infty = \underset{t\in \icc{a}{b}}{\sup}\{|f(t)|\}\)
}
Graphiquement, la boule ouverte de centre la fonction \(\sin\) et de rayon \(1\) pour la norme infini est alors:\<

\begin{center}
   \begin{tikzpicture}[
      >=stealth,scale=1,line cap=round,
      bullet/.style={circle,inner sep=0.5pt,fill}
   ]
      \draw[->] (0,-2) -- (0, 3);
      \draw[] (-0.1,0) -- (0.1, 0) node at (-0.3, 0) {$0$};
      \draw[] (-0.1,1) -- (0.1, 1) node at (-0.3, 1) {$1$};
      \draw[] (-0.1,2) -- (0.1, 2) node at (-0.3, 2) {$2$};
      \draw[] (-0.1,-2) -- (0.1, -2) node at (-0.45, -2) {$-2$};
      \draw[] (-0.1,-1) -- (0.1, -1) node at (-0.45, -1) {$-1$};

      \draw[color=BrightRed1, domain=0:12, smooth]   plot (\x,{sin(\x r)}) node[right] {$\sin(t)$}; 
      \draw[color=BrightRed1!60,name path=A, domain=0:12, smooth, dashed]   plot (\x,{sin(\x r) + 1}); 
      \draw[color=BrightRed1!60,name path=B, domain=0:12, smooth, dashed]   plot (\x,{sin(\x r) - 1}); 
   \end{tikzpicture}   
\end{center}

\subsection*{\subsecstyle{Exemples de boules exotiques {:}}}
Si on considère maintenant l'ensemble des cases d'un échiquier \(6 \times 6\) et qu'on définit la distance entre deux cases comme étant le nombre de coup qu'un cavalier doit effectuer pour arriver à la seconde en partant de la première, on définit alors un espace métrique et la boule de centre le centre d'un plateau et de rayon 1 est représentée par:
\begin{center}
   \begin{tikzpicture}[scale=0.5, every node/.style={minimum size=.5cm-\pgflinewidth, outer sep=0pt}]
      \node[fill=BrightBlue1] at (0.5,1.5) {};      
      \node[fill=BrightBlue1] at (1.5,0.5) {};      
      
      \node[fill=BrightBlue1] at (0.5,3.5) {};      
      \node[fill=BrightBlue1] at (1.5,4.5) {};      
      
      \node[fill=BrightBlue1] at (4.5,1.5) {};      
      \node[fill=BrightBlue1] at (3.5,0.5) {};

      \node[fill=BrightBlue1] at (4.5,3.5) {};      
      \node[fill=BrightBlue1] at (3.5,4.5) {};

      \node[fill=BrightRed1] at (2.5,2.5) {};
      \draw[step=1cm,color=black, line width = 0.2mm] (0,0) grid (5,5);
   \end{tikzpicture}
\end{center}



\chapter*{\chapterstyle{II --- Continuité}}
\addcontentsline{toc}{section}{Continuité}

\chapter*{\chapterstyle{II --- Connexité}}
\addcontentsline{toc}{section}{Connexité}

\chapter*{\chapterstyle{II --- Compacité}}
\addcontentsline{toc}{section}{Compacité}



